\sloppy
\chapter{Related Work}
\addtocontents{toc}{\protect\setcounter{tocdepth}{2}}
\section{Introduction}

In \citep{Kaliyar2021FakeBERT}, a BERT-based deep learning model called FakeBERT was proposed specifically for fake news detection on social media platforms, integrating parallel CNN layers into the BERT architecture for multi-scale feature extraction over the LIAR and ISOT datasets. The model achieved 98.90\% accuracy, effectively capturing contextual semantic relationships in text and substantially outperforming traditional ML models. However, it is computationally expensive, requires large GPU memory, and is limited to English-language content. The hybrid model proposed in this thesis pursues a similar transformer-centred idea, but layers a BiLSTM instead of parallel CNNs on top of BERT's contextual embeddings, and is evaluated across multiple datasets rather than a single one.

Multiple classical machine learning classifiers were systematically compared for fake news detection in \citep{Aphiwongsophon2021Detecting}, using linguistic and statistical text features extracted via TF-IDF with Logistic Regression, SVM, Naive Bayes, and Random Forest. These models are simple to implement, computationally cheap, and interpretable, making them suitable for resource-constrained environments. Their accuracy, however, trails deep learning approaches because they cannot capture semantic context or word order in text. The system proposed in this thesis instead relies on BERT's contextual embeddings, trading some of that simplicity for a considerably richer representation of meaning.

A hybrid deep learning architecture combining a CNN for local feature extraction and an RNN for sequential pattern learning was proposed in \citep{Nasir2021FakeNews}, trained with GloVe word embeddings on the ISOT and FA-KES datasets. It captures both local and sequential text features simultaneously and generalizes better than single-model approaches, though it requires significant training data, has slow inference time, and struggles with short text such as tweets and headlines. Our proposed BERT + BiLSTM model addresses the short-text weakness by relying on BERT's pre-trained contextual representations rather than embeddings learned from scratch.

Stylometric and linguistic style features that distinguish fake news from credible news, such as sentence length, punctuation patterns, and readability scores, were analysed with an SVM classifier in \citep{Przybyla2020Capturing}. This approach does not require an external knowledge base, works purely on writing style, and is lightweight and fast at inference time. Its main weakness is that style-based features are easily mimicked by a sophisticated fake news creator, giving it low robustness against adversarial content, a concern that contextual, semantics-driven models such as the one proposed here are comparatively less susceptible to.

An evidence-aware fake news detection approach was proposed in \citep{Popat2019DeClarE}, retrieving external web evidence and combining it with claim credibility assessment using an attention-based LSTM over the Snopes and PolitiFact datasets. It incorporates external knowledge for verification and its attention mechanism provides interpretability, performing well on political claims. It is, however, dependent on the availability of external evidence, slow due to web retrieval, and fails when evidence is unavailable or biased. The present work does not perform evidence retrieval, focusing instead on the classification model itself, though Chapter~5 identifies retrieval-augmented evidence grounding as a natural future extension.

A comprehensive fake news benchmark dataset, FakeNewsNet, was developed in \citep{Shu2020FakeNewsNet}, incorporating news content, user social engagement, and spatiotemporal propagation features drawn from PolitiFact and GossipCop. It is one of the most comprehensive publicly available datasets and enables multi-modal detection research, though its collection is time-consuming and its social-context features require platform API access that is increasingly restricted. This thesis uses the more text-centric LIAR and ISOT datasets, which are better suited to a controlled, text-only model comparison.

State-of-the-art deep learning techniques for fake news detection, covering transformer models, graph neural networks, and multi-modal approaches, were surveyed in \citep{Mehta2022RecentTrends}. The survey offers comprehensive coverage of recent methods and identifies research gaps, though it proposes no novel model and its coverage is limited to English-language studies. The systematic sequential-versus-transformer-versus-hybrid comparison carried out in this thesis addresses one of the concrete gaps this kind of survey highlights.

BERT contextual embeddings were combined with an enhanced LSTM architecture and an additional attention mechanism in \citep{Rai2022FakeNews}, trained on the LIAR and ISOT datasets. This hybrid achieved superior performance over standalone BERT or LSTM, with the attention layer improving interpretability and robustness to text length variation. Its drawbacks are high model complexity, longer training time, and the need for substantial labelled data for effective fine-tuning. It is the closest published work in spirit to the BERT + BiLSTM model proposed here, though this thesis omits the additional attention layer to isolate the specific contribution of sequential modelling on top of BERT.

Psychological and cognitive bias patterns embedded in fake news content were investigated in \citep{CarrascoFarre2022Fingerprints} using computational linguistics and psycholinguistic feature extraction with the LIWC tool alongside ML classifiers. This offers a novel perspective combining psychology and NLP, with interpretable features applicable across domains and languages, but it is limited by small dataset size and psycholinguistic tools that may not generalize across cultural and linguistic contexts. This dimension of analysis is complementary to, and outside the scope of, the purely text-based transformer models compared in this thesis.

An optimized ensemble learning model using a metaheuristic harmony search algorithm for automatic feature selection and classifier weight optimization was proposed in \citep{Huang2020FakeNews}, combining SVM, Random Forest, and Logistic Regression across multiple news datasets. Automatic feature selection reduces dimensionality and the ensemble approach improves robustness, outperforming individual classifiers, though at the cost of high computational overhead due to the metaheuristic optimization step and an implementation too complex for real-time detection. The transformer-based models evaluated in this thesis avoid this optimization overhead entirely by relying on pre-trained representations.

A foundational data mining perspective on fake news detection on social media was provided in \citep{Shu2017FakeNewsDetection}, proposing a taxonomy of detection methods spanning knowledge-based, style-based, propagation-based, and credibility-based approaches. It offers a comprehensive taxonomy covering diverse detection paradigms and remains a foundational reference for the research community, although it is a survey with no experimental results and limited coverage of deep learning and transformer-based models. This thesis sits within the credibility-based, content-driven branch of that taxonomy.

A comprehensive survey of fake news fundamental theories, detection methods, and future opportunities was presented in \citep{Zhou2020Survey}, reviewing knowledge-based, style-based, propagation-based, and credibility-based detection methods across multiple benchmark datasets. It provides thorough coverage of detection strategies and identifies open problems, but proposes no novel detection model and gives limited coverage to multimodal and LLM-based approaches, gaps that motivate the future work outlined in Chapter~5 of this thesis.

CSI (Capture, Score, Integrate), a hybrid deep learning model combining text content features with user behaviour and propagation patterns, was proposed in \citep{Ruchansky2017CSI} using a hybrid RNN architecture evaluated on a CIKM benchmark dataset. It was the first model to integrate content, user, and propagation signals, significantly outperforming text-only baselines on social media platforms. It requires social engagement data that may be unavailable, has a complex multi-module architecture, and generalizes poorly to non-social-media news. The present work deliberately stays text-only, trading the additional signal CSI exploits for a simpler, more broadly applicable model.

A landmark empirical study on the spread of true and false news online across Twitter was conducted in \citep{Vosoughi2018Spread}, analysing 126,000 news stories tweeted by three million users over ten years and finding that false news spreads roughly six times faster than true news and reaches more people. It is the largest empirical study of its kind and provides strong motivation for automated fake news detection, although it is purely observational, proposes no detection model, and is limited to Twitter platform data. This finding is one of the central motivations cited in Chapter~1 of this thesis.

The LIAR benchmark dataset, containing 12,836 manually labelled short statements from PolitiFact with six fine-grained truthfulness labels, was introduced in \citep{Wang2017Liar} together with metadata such as speaker, subject, context, and party affiliation, and baseline CNN and BiLSTM classifiers. It was the first large-scale manually labelled fake news dataset and its metadata-rich annotations support multi-feature modelling, though it is limited to the political speech domain and its fine-grained labels make classification harder. LIAR is one of the two benchmark datasets used to evaluate the models compared in this thesis.

LLaMA 2, a family of open-source large language models ranging from 7B to 70B parameters, was introduced in \citep{Touvron2023Llama2}, fine-tuned for dialogue and instruction following and evaluated against GPT-4 and other proprietary models. It is freely available for research, performs competitively against closed models, and enables privacy-preserving local deployment, although it requires significant GPU resources, may still produce hallucinations, and has more limited multilingual capability than GPT-4. Fine-tuning a model such as LLaMA 2 on fake news data is identified as a future direction in Chapter~5.

Retrieval-Augmented Generation (RAG), a framework combining the parametric knowledge of pre-trained LLMs with non-parametric retrieval from external knowledge bases, was proposed in \citep{Lewis2020RAG}, combining dense passage retrieval with a sequence-to-sequence generator and evaluated on open-domain QA, fact verification, and knowledge-grounded dialogue. RAG enables LLMs to access up-to-date external evidence and reduces hallucination without full model retraining, but its retrieval quality depends heavily on the coverage of the knowledge base and it introduces additional inference latency. Grounding fake news classification decisions in evidence retrieved by a RAG-style pipeline is proposed as future work in Chapter~5.

LIME (Local Interpretable Model-agnostic Explanations), a technique for explaining the predictions of any black-box classifier by approximating it locally with an interpretable model, was proposed in \citep{Ribeiro2016LIME}. It works with any classifier regardless of architecture and provides human-readable explanations that enable trust and debugging, though its local explanations may not reflect global model behaviour and it can be computationally expensive for large input spaces. Applying LIME to explain which words drove a BERT-based REAL or FAKE prediction is one of the explainability directions discussed in Chapter~5.

SHAP (SHapley Additive exPlanations), a unified framework for interpreting model predictions using Shapley values from cooperative game theory, was proposed in \citep{Lundberg2017SHAP}. It is theoretically grounded, consistent, and locally accurate, enabling both global feature importance and individual prediction explanations, but it is computationally expensive for large feature spaces and deep models and requires careful interpretation. Combining SHAP with attention maps to quantify word-level contributions in the BERT-based classifier is likewise identified as future work in Chapter~5.

\section{Summary}

The literature reviewed in this chapter spans classical machine learning classifiers, hybrid deep learning architectures, evidence-aware and psycholinguistic approaches, benchmark dataset papers, broad surveys, and recent large language model and explainability techniques. Across this body of work, a common pattern emerges: classical and purely statistical methods remain cheap and interpretable but cannot capture semantic context, while more capable deep learning and transformer-based methods trade that simplicity for higher accuracy at greater computational cost, and few of the reviewed works directly compare sequential, transformer-based, and hybrid architectures under the same controlled conditions. We believe that a systematic, apples-to-apples comparison of a standalone BiLSTM, a standalone BERT model, and a hybrid BERT + BiLSTM model, evaluated consistently across multiple benchmark datasets, addresses this gap and provides clearer practical guidance on when the added complexity of a hybrid architecture is actually justified.
